
\begin{wrapfigure}{r}{0.45\textwidth}
\vspace{-5mm}
\includegraphics[width=0.45\textwidth]{figures/GBM_ULM_Image_Timeline_Published.png}
\caption{\textbf{ Volumetric 3D transcranial super-resolution ultrasound imaging} was used non-invasively in an \textit{in vivo} longitudinal study of glioblastoma in rodents. A significant ($p<0.05$) reduction in the appearance of functional vasculature was observed over three weeks along with reduced vascular flow peripheral to the tumor. This demonstrates that super-resolution imaging can acquire high-quality transcranial vascular images non-invasively, establishing the maturity level required for microvascular characterization of neurological disease including stroke.\label{fig:gbm} (From~\cite{mccall2023longitudinal}) }
\vspace{-3mm}
\end{wrapfigure}

\subsection{Significance of portable, low-cost neurovascular imaging and high-precision interventions in stroke:}
With nearly 800,000 cases annually in the U.S., ischemic stroke remains governed by a narrow therapeutic window in which every minute of delayed reperfusion results in the loss of approximately 1.9 million neurons. The two principal interventions---intravenous tPA (effective within 4.5~hours) and endovascular thrombectomy (EVT, for large-vessel occlusions at specialized centers)---leave roughly 500,000 patients per year untreated due to time-window expiration, tPA contraindications, vessel inaccessibility to catheter-based approaches, or lack of proximity to a comprehensive stroke center. Efforts to compress the time from symptom onset to treatment have driven the development of mobile stroke units equipped with onboard CT and telemedicine, as well as regionalized stroke triage protocols that route patients directly to thrombectomy-capable centers. While these initiatives have improved door-to-needle and door-to-groin times, they remain constrained by the cost and infrastructure requirements of CT/MRI-based diagnosis and the limited geographic availability of neurointerventional expertise.

This gap motivates the development of a portable, low-cost imaging platform that can be deployed at the point of care---in ambulances, emergency departments, and community hospitals---to both diagnose and initiate treatment of ischemic stroke. Ultrasound is the only modality that is simultaneously portable, real-time, free of ionizing radiation, and capable of delivering therapeutic energy. Recent advances in functional ultrasound imaging (fUSI) and super-resolution ultrasound have demonstrated sensitivity to hemodynamic changes and microvascular architecture deep within the brain that approach or exceed the resolution of conventional neuroimaging~\cite{demene2017functional,rabut20194d,mccall2023non}. Critically, because diagnostic and therapeutic ultrasound share the same physical platform, imaging and intervention can be inherently co-registered and performed concurrently---enabling a theranostic system in which high-resolution neurovascular assessment directly guides targeted sonothrombolysis without the delays, costs, or logistical constraints of separate imaging and interventional suites.

\subsection{Significance of super-resolution imaging for vascular and microvascular characterization of stroke and therapeutic interventions:}
Super-resolution ultrasound imaging (also called ultrasound localization microscopy) is a rapidly emerging technology for \textit{in vivo} microvascular imaging being developed by several groups worldwide, including ours~\cite{christensen2020super}. By localizing the centers of freely circulating microbubble contrast agents, it resolves microvascular structures far beyond the acoustic diffraction limit. This technique has been used to image the microvasculature of the rodent brain with scalp removal~\cite{renaudin2022functional}, through a thinned skull~\cite{errico2015ultrafast}, with a craniectomy~\cite{milecki2021deep}, in several other organs~\cite{christensen2014vivo,foiret2017ultrasound,song2017improved,huang2020short,lin20173,lowerison2020ultrasound}, and in humans~\cite{demene2021transcranial}.

Volumetric 3D super-resolution is essential because 2D imaging, even with mechanical scanning, cannot super-resolve transverse vessels, underestimates velocity and flow volume, and requires impractically long scan times~\cite{chavignon20213d}. Volumetric super-resolution has been demonstrated \textit{in vivo} in transcranial rat and mouse brain with resolutions of 6.5--31~$\mu$m and flows between 2--100~mm/s~\cite{chavignon20213d,demeulenaere2022vivo,renaudin2022functional}, and completely non-invasively by us~\cite{mccall2023non}. We recently demonstrated the power of this approach in a longitudinal 3D visualization of microvascular disruption and perfusion changes during glioblastoma evolution (Fig.~\ref{fig:gbm})~\cite{mccall2023longitudinal}, establishing that the technology has reached the maturity level for sophisticated microvascular characterization of neurological disease.

For stroke specifically, super-resolution imaging offers transformative diagnostic and therapeutic significance. In ischemic stroke, the ability to resolve individual microvessels and quantify their flow enables: (1)~direct visualization of vascular occlusion sites including in small vessels inaccessible to conventional angiography; (2)~quantitative mapping of the penumbral region where reduced but salvageable perfusion exists, which is critical for treatment decisions; and (3)~real-time monitoring of microvascular reperfusion during and after therapeutic intervention. This last capability is particularly significant because it provides direct, quantitative feedback on the efficacy of sonothrombolysis at the microvascular level---information that is currently unavailable with any existing clinical imaging modality. The ability to visualize restoration of flow in individual microvessels during therapy would enable adaptive, precision-guided treatment that maximizes therapeutic benefit while minimizing risk.

\subsection{Significance of ultrasound as a non-invasive targeted sonothrombolysis tool:}
Ultrasound-mediated thrombolysis has demonstrated clinical efficacy---the CLOTBUST trial showed that continuous 2~MHz transcranial Doppler insonation combined with IV tPA increased complete recanalization rates compared to tPA alone~\cite{alexandrov2004ultrasound}, and subsequent studies explored sonothrombolysis in patients ineligible for standard thrombolysis~\cite{eggers2005sonothrombolysis,eggers2008sonothrombolysis}. The underlying mechanisms are well established: acoustic streaming and radiation forces enhance transport of thrombolytic agents into the clot, while cavitation generates microstreaming and mechanical stresses that directly disrupt the fibrin matrix~\cite{datta2006correlation,everbach2000cavitational}. These mechanical effects can be produced without thrombolytic drugs, opening the possibility of purely mechanical, non-pharmacological clot dissolution that avoids the systemic bleeding risks of tPA. Our laboratory has extensive experience in sonothrombolysis, including development of intravascular forward-looking transducers for microbubble-mediated sonothrombolysis~\cite{kim2017intravascular,kim2021multi,goel2021safety}, demonstration that low-boiling-point nanodroplets achieve superior clot dissolution compared to microbubbles for both aged and retracted clots~\cite{kim2020comparison,kim2022analysis,goel2021nanodroplet}, combination strategies using magneto-sonothrombolysis~\cite{zhang2021magneto} and low-dose tPA with microbubbles~\cite{goel2020examining}, and a comprehensive review of sub-micron cavitation-enhancing agents~\cite{bautista2023current}.

However, a fundamental limitation of all prior clinical sonothrombolysis approaches is the absence of targeting. The CLOTBUST and related trials used unfocused transcranial Doppler transducers that insonified hundreds of cubic centimeters of brain parenchyma with no ability to confine energy to the clot. The consequences of this were demonstrated starkly in the TRUMBI trial, where low-frequency (300~kHz) unfocused insonation combined with tPA produced unacceptably high rates of symptomatic intracranial hemorrhage, leading to early termination~\cite{daffertshofer2005transcranial}. The TRUMBI outcome illustrates a central problem: when ultrasound energy is distributed over large tissue volumes rather than confined to the thrombus, the risk of off-target bioeffects---hemorrhagic transformation, edema, and unintended cavitation in healthy parenchyma---rises dramatically.

Our approach represents a fundamentally different paradigm: \textit{targeted} sonothrombolysis, in which focused therapeutic ultrasound is delivered exclusively to the clot volume. Using image-guided focusing, the therapeutic beam is confined to a focal volume of a few~mm$^3$---orders of magnitude smaller than the hundreds of cm$^3$ insonified by unfocused approaches. Super-resolution imaging localizes the thrombus with micron-scale accuracy and the co-registered therapeutic beam delivers sonothrombolysis with 750~$\mu$m focal diameter and 50~$\mu$m positional accuracy, ensuring that acoustic energy is deposited at the clot and not in surrounding tissue. Real-time imaging feedback enables continuous monitoring of hemodynamic restoration during treatment, allowing the acoustic dose to be adapted based on observed reperfusion. This precision targeting directly addresses the safety concerns that have limited clinical translation of sonothrombolysis: by restricting mechanical and cavitational effects to the thrombus itself, the risk of hemorrhagic transformation that plagued untargeted trials is substantially mitigated while therapeutic efficacy at the clot is maximized.

\subsection{Significance of perfluorocarbon contrast agents as dual-purpose imaging and sonothrombolytic agents:}
Perfluorocarbon-based ultrasound contrast agents occupy a uniquely dual role in the proposed theranostic platform: they serve simultaneously as the basis for super-resolution imaging and as potent enhancers of sonothrombolysis. This agent family encompasses two complementary formulations: conventional gas-filled microbubbles (1--3~$\mu$m) and phase-change nanodroplets ($\sim$200~nm liquid perfluorocarbon particles that convert to gas-filled microbubbles upon acoustic activation)~\cite{bautista2023current}. Both share the same core perfluorocarbon chemistry and lipid shell composition, but their distinct sizes and physical states confer complementary mechanisms for both imaging and therapy.

As imaging agents, microbubbles are well established for super-resolution ultrasound---their point-like acoustic signatures enable sub-diffraction localization, and their nonlinear scattering properties provide the high contrast-to-tissue ratio essential for reliable detection and tracking~\cite{christensen2020super,gessner2010high,gessner2013acoustic}. Our group has pioneered super-harmonic imaging, which exploits the broadband nonlinear emissions of microbubbles using dual-frequency transducers to achieve outstanding contrast-to-tissue ratios and enable visualization of microvascular structures on the order of 100~microns~\cite{gessner2013acoustic,gessner2010high,gessner2012mapping,rao2015fingerprint,shelton2015quantification,shelton2020microvascular}. Phase-change nanodroplets, once activated, produce microbubbles with similar nonlinear properties, and their smaller pre-activation size may allow detection of agent accumulation within thrombi that are inaccessible to circulating microbubbles. The essential absence of background tissue clutter in super-harmonic imaging enables reliable segmentation and quantitative microvascular analysis.

As therapeutic agents, microbubbles enhance sonothrombolysis through stable and inertial cavitation, generating localized mechanical forces---microstreaming, microjets, and shear stresses---that erode the fibrin clot matrix from the surface~\cite{datta2006correlation,everbach2000cavitational}. However, conventional microbubbles are confined to the vascular lumen and can only act on the external clot surface, limiting their efficacy against dense, retracted thrombi. Phase-change nanodroplets overcome this fundamental limitation: their sub-micron size allows them to diffuse into the fibrin matrix, where ultrasound-triggered phase change generates microbubbles that cavitate from within the clot interior, producing a novel ``permeation and internal disruption'' mechanism~\cite{kim2020comparison,goel2021nanodroplet}. Our studies have demonstrated that this approach achieves clot dissolution in dense, aged clots that is unachievable with conventional microbubbles alone~\cite{kim2022analysis,kim2020comparison}. Critically, we have shown that the combination of conventional microbubbles and phase-change nanodroplets achieves optimal clot mass reduction that exceeds either agent alone~\cite{kim2022analysis}---external cavitation from microbubbles and internal cavitation from activated nanodroplets act synergistically. Furthermore, dual-frequency ultrasound excitation significantly enhances thrombolysis efficacy compared to single-frequency approaches~\cite{wu2021dual}, a strategy that aligns naturally with the dual-frequency array architecture of our proposed platform.

The dual-purpose use of these agents in our platform creates a synergistic feedback loop: the same contrast agents used to image the vasculature and localize the thrombus are the agents that enhance clot dissolution when therapeutic ultrasound is applied. As agents circulate through partially recanalized vessels during treatment, they simultaneously provide real-time imaging feedback on the degree of flow restoration, enabling the operator to assess therapeutic efficacy and adjust the treatment in real time. This tight integration of imaging and therapy through a single agent family eliminates the need for separate diagnostic and therapeutic administrations and ensures inherent co-registration between the diagnostic and therapeutic functions of the system. Furthermore, advances in super-harmonic imaging with dual-frequency arrays~\cite{kierski2020superharmonic,kierski2022acoustic} allow separation of contrast agent signal from tissue without reliance on spatiotemporal filtering, enabling detection of both freely circulating and stationary agents---a capability essential for monitoring clot-associated nanodroplets and microbubbles.

\subsection{Significance of the research team:} Dr. Pinton (UNC) has extensive experience in ultrasound imaging, super-resolution imaging, transcranial ultrasound, and numerical wave propagation. Dr. Dayton (UNC), recipient of the 2020 IEEE UFFC Hertz Ultrasonics Award, has over 20 years of experience and over 200 publications in the field of microvascular imaging and contrast ultrasound-based technologies. Dr. Lee, MD, Ph.D., is vice chair of translational research at UNC, has worked with Drs. Pinton and Dayton on prior clinical trials, and will provide clinical experience in evaluating the translational relevance of the proposed imaging and therapeutic approaches for stroke.
